% SOURCE_AUTHORITY: sga5_fr_workpass.tex, lines 14542--14573 (LF-normalized unit).
% SOURCE_SHA256: 791F4EFFC5E02832D5D77ED03518C8156D6F07E4C8238B03545DB93D883FBB28
\subsection*{§ 2. Correspondencia de Frobenius}

\subsubsection*{n\textsuperscript{o} 1. Definición}

Sean $X$ un preesquema de característica $p$ y $F$ un haz de conjuntos sobre $X_{\et}$. Para todo $X$-preesquema étale $U$ se tiene
\[
(\fr_X)_*(F)(U)=F(\fr_X^{-1}(U))=F(U^{(p/X)}),
\]
y el morfismo de Frobenius $\Fr_{U/X}:U\to U^{(p/X)}$ define una aplicación
\[
F(\Fr_{U/X}):\fr_{X*}(F)(U)\longrightarrow F(U),
\]
manifiestamente funtorial en $U$ (pues $\Fr_{U/X}$ es funtorial en $U$). Se obtiene así un morfismo de haces sobre $X_{\et}$:
\[
(\fr_X)_*(F)\longrightarrow F.
\]
Ahora bien, cuando $U$ es étale sobre $X$, $\Fr_{U/X}$ es un isomorfismo (§ 1 n$^\circ$2, prop. 2c)), luego lo mismo ocurre con $F(\Fr_{U/X})$; resulta que el morfismo precedente es un isomorfismo. En el caso en que $F$ esté provisto de una estructura algebraica (por ejemplo, si es un haz de $\Lambda$-módulos, siendo $\Lambda$ un anillo base dado), este isomorfismo es compatible con la estructura algebraica. La inversa
\[
F(\Fr_{/X})^{-1}:F\longrightarrow \fr_{X*}(F)
\]
de este isomorfismo define, puesto que $(\fr_X)^*$ es adjunto por la izquierda del funtor $(\fr_X)_*$, un morfismo
\[
\Fr^*_{F/X}:(\fr_X)^*(F)\longrightarrow F.
\]
Se puede decir incluso algo más, pues, al ser $\fr_X$ un homeomorfismo universal (§ 1 n$^\circ$2, prop. 2a)), los funtores $(\fr_X)^*$ y $(\fr_X)_*$ son cuasi-inversos uno del otro y, por consiguiente, $\Fr^*_{F/X}$ es un isomorfismo, al igual que el morfismo $F(\Fr_{/X})^{-1}$ del que procede.

\begin{statement}{Definición 1}
Dados un preesquema $X$ de característica $p$ y un haz de conjuntos $F$ sobre $X_{\et}$, se llama correspondencia de Frobenius sobre $(X,F)$ al par $(\fr_X,\Fr^*_{F/X})$ formado por el endomorfismo de Frobenius $\fr_X$ de $X$ y el isomorfismo
\[
\Fr^*_{F/X}=(F(\Fr_{/X})^{-1})^\sharp:(\fr_X)^*(F)\longrightarrow F.
\]
\end{statement}
